\newpage
\subsection{Temperaturmessung mit PT-100}\label{sec:1-5}
\subsubsection{Aufgabenstellung}

Eine Wheatstone Brücke soll so dimensioniert werden, dass mit einem PT-100
Metall-Widerstandsthermometer die Raumtemperatur gemessen werden kann. 
Die Brücke soll bei einer Temperatur von \SI{0}{\degreeCelsius} die
Abgleichbedingung erfüllen und dort auch das empfindlichste
Widerstandsverhältnis $R_1/R_2$ aufweisen. Mittels der Hilfsspannung $U_H$ soll
eingestellt werden, dass bei einer Temperatur von \SI{100}{\degreeCelsius} die
Brückenspannung $U_0=\SI{50}{\milli\volt}$ beträgt. Dabei muss beachtet werden,
dass über den PT-100 nicht mehr als \SI{4}{\milli\ampere} fließen, um
Eigenerwärmung zu vermeiden.

\subsubsection{Verwendetes Material}

\begin{table}[h]
	\centering
	\caption{Verwendetes Material}
	\begin{tabularx}{\textwidth}{XXX}
		\textbf{Material} & \textbf{Verwendung} & \textbf{Anmerkungen}\\ \midrule
		HM8001 & Hilfsspannung ($U_H$) & \\
		M-4600 & Nullinstrument ($U_0$) & Bezeichung: NI\\
		PT-100 Simulator & Zur Dimensionierung ($R_1$) & Modul auf der Messbrücke\\
		PT-100 & Temperatursensor ($R_1$) & \\
	\end{tabularx}	
\end{table}

\subsubsection{Versuchsaufbau}

Zuerst wird der PT-100 Simulator mit dem Steckplatz $R_1$ verbunden. An diesem
lassen sich Widerstandswerte eines PT-100 zu diskreten Temperaturen zwischen
\SI{0}{\degreeCelsius} und \SI{100}{\degreeCelsius} einstellen. Nach der
Dimensionierung in \autoref{sec:dim} der Brücke kann der PT-100 eingesteckt
werden, um die Temperatur zu messen.

\begin{figure}[h]
	\centering
	\begin{circuitikz}[scale=0.9, transform shape]
		\draw (0,4) to[open, v=$U_H$, o-o] (0,0)
			-- ++(2, 0) coordinate (k1)
			to[R=$R_2$, *-] ++(0,2) coordinate (k2)
			to[thRp=$R_1$, -*] ++(0,2) coordinate (k3)
			-- ++(-2,0);
		\draw (k1) -- ++(3,0)
			to[vR=$R_4$] ++(0,2) coordinate (k4)
			to[R=$R_3$] ++(0,2) -- ++(-3,0);
		\draw (k2) to[ai, v^=$U_0$, *-*] (k4);
		\node[yshift=-20pt] at ($(k2)!0.5!(k4)$) {$R_0$};
	\end{circuitikz}
	\caption{Messschaltung}
\end{figure}

\subsubsection{Dimensionierung}\label{sec:dim}

Bei \SI{0}{\degreeCelsius} weist der PT-100 Simulator \SI{100}{\ohm} auf. Mit $R_4$
im Bereich \SI{120}{\ohm} lassen sich aus der Abgleichbedingung die Widerstände
$R_2$ und $R_3$ geeignet wählen.
\[
R_1 R_4 = R_2 R_3 = \SI{100}{\ohm} \cdot \SI{120}{\ohm}
	\implies R_2 = \SI{100}{\ohm}, R_3 = \SI{120}{\ohm}
\]
Mit dieser Wahl gilt $R_1:R_2 = 1:1$ für maximale Empfindlichkeit. Nun
kann eine geeignete Hilfsspannung eingestellt werden. Der Querstrom im zweig
des PT-100 ist
\[
I_{R_1} = \frac{U_H}{R_1 + R_2} < \SI{4}{\milli\ampere}
	\implies U_H < \SI{4}{\milli\ampere} \cdot \SI{200}{\ohm} = \SI{0.8}{\volt}
\]
Beim ermitteln der zweiten Randbedingung, dass $T=\SI{100}{\degreeCelsius}
\mapsto U_0=\SI{50}{\milli\volt}$ wurde bei $U_H=\SI{0.8}{\volt}$ bereits
$\SI{-58.21}{\milli\volt}$ festgestellt und beschlossen mit diesem Wert
fortzufahren. Die Auswirkung dieser Abweichung der Anforderungen werden im
Anschluss diskutiert.

\subsubsection{Messung}

Mit der dimensionierten Brücke wird nun der PT-100 Eingesteckt und eine
Spannung von $U_{0,\text{RT}} = \SI{-16.76}{\milli\volt}$ festgestellt.

\subsubsection{Auswertung}

Mit der Gleichung \cite[S.~41, Gl.~2.5]{EMTPraktikum2025} für die Brückenspannung
\begin{equation}
U_0 \approx U_H \frac{-\frac{\Delta R_1}{R_1}}{\left(1+\frac{R_2}{R_1}\right)\left(1+\frac{R_3}{R_4}\right)}\label{eq:U0temp}
\end{equation}
lässt sich auf den tatsächlichen Widerstandswert des PT-100 bei dieser Temperatur zurückschließen. Umformen auf $\Delta R_1$ 
\[
\Delta R_1 = -\frac{U_{0,\text{RT}}}{U_H}(R_1+R_2)\left(1+\frac{R_3}{R_4}\right) = \SI{8.38}{\ohm}
\]

Mit der linearen Taylorapproximation für Temperaturwiderstände
\cite[S.~17, Gl.~1.24]{EMTPraktikum2025} kann dann die Temperatur berechnet
werden. Für den PT-100 beträgt der Temperaturkoeffizient erster Ordnung
$\alpha=\SI{3.85e-3}{\per\kelvin}$.
\[
\Delta R_1(\theta) + R_1 = R_1(1+\alpha\theta) \implies \theta = \frac{1}{\alpha}\frac{\Delta R_1}{R_1} = \SI{21.77}{\degreeCelsius}
\]
Ein Referenzthermometer liefert einen Anzeigewert von ca. \SI{21}{\degreeCelsius}.

\subsubsection{Erkenntnisse}

Die berechnete Temperatur liegt sehr nahe an dem gemessenen Wert. Die
Entscheidung $U_H$ auf \SI{0.8}{\volt} zu lassen, hatte keine negativen
Auswirkungen auf das Ergebnis, da sich für eine größere Hilfsspannung $U_0$
proportional erhöht (siehe \autoref{eq:U0temp}) und daher das Verhältnis gleich
bleibt.